\documentclass[../numerical_methods.tex]{subfiles}
\begin{document}
\section{Aula Extra}
\subsection{Motivações}
\begin{itemize}
	\item Noções de Erros e Pontos Flutuantes.
\end{itemize}
\subsection{Noções de Erros e Pontos Flutuante}
A ideia de pontos flutuantes é encontrar uma forma de representar os números reais, essencialmente infinitos tanto em tamanho quanto em granularidade, em um computador, que não tem processamento e nem memória infinita.
\begin{def*}
	O conjunto \(\mathbb{F} = \{\text{representação de x}:\; x\in \mathbb{R}\}\) é o conjunto cujos elementos \(\overline{x}\in \mathbb{F}\) são chamados de \textbf{pontos flutuantes}. \(\square\)
\end{def*}
A representação deles consiste primeiro em entender bases numéricas para inteiros, como a já conhecida base decimal e a base binária. Formalmente,
\begin{def*}
	Dado um inteiro N, sua versão na base \(\beta \) é dada por
	\[
		N = d_{k}\beta^{k}+d_{k-1}\beta^{k-1} +\dotsc +d_1\beta^{1} + d_{0}\beta^{0},
	\]
	onde \(0\leq d_{i}\leq (\beta-1)\) para \(i=0, \dotsc , k\). Denotamos \textbf{N na base \(\mathbf{\beta }\)} por \(N = (d_kd_{k-1}\dotsc d_1d_0)_\beta. \; \square\)
\end{def*}
\begin{example}[Conversão da Base Binária para Decimal]
	Vamos converter o número \((10011)_{2}\) para a base decimal:
	\[
		(10011)_{2} = 1 \times 2^{4} + 0 \times 2^{3} + 0 \times 2^{2} + 1 \times 2 + 1\times 2^{0} = 19 = 1\times 10^{1} + 9\times 10^{}0 = (19)_{10}.
	\]


	Para o caminho contrário, dado N um número inteiro na base decimal, desejamos representar N na forma
	\[
		(d_k d_{k-1}\dotsc d_1d_0)_{2} = d_{k}2^{k}+\dotsc +d_1 2^{1} + d_{0}2^{0};
	\]
	logo,
	\begin{align*}
		N & = 2 (d_k \times 2^{k-1} + d_{k-1}2^{k-1}+\dotsc +d_2 2^{1} + d_1) + d_{0} \\
		  & = 2 (2 (d_{k} 2^{k-2} + d_{k-1}2^{k-3}+\dotsc +d_2)+d_1)+d_0              \\
		  & \vdots
	\end{align*}
	Assim adiante, até o quociente ficar nulo. Tente mostrar que \(2(11101)_{2} = (29)_{10}.\)
\end{example}
Com base no exemplo acima, dá pra montar um algoritmo bem simples que converte da decimal pra binária:
\begin{flushleft}

	Entrada: \(N \in \mathbb{Z}\)

	Saída: \(N = (d_kd_{k-1}\ldots d_1d_0)_2\)

	\vspace{0.4cm}

	\(k = 0;\)

	Calcule \(R\) e \(Q\) tais que \(N = 2Q + R;\)

	\(d_k = R;\)

	\vspace{0.3cm}

	Enquanto \(Q \neq 0\) faça

	\hspace{1cm} \(k = k + 1;\)

	\hspace{1cm} \(N = Q;\)

	\hspace{1cm} Calcule \(R\) e \(Q\) tais que \(N = 2Q + R;\)

	\hspace{1cm} \(d_k = R;\)

	Fim do Enquanto

\end{flushleft}

Agora que temos a noção de como funciona pros inteiros, podemos explicar o raciocínio para um número N no intervalo \((0, 1):\) se \(N=0, d_1d_2 \dotsc d_{k}\), nosso objetivo é escrever
\[
	N = (0, d_1d_2\cdots d_k)_{2} = d_1 2^{-1} + d_2 2^{-2} + \dotsc + d_{k}2^{-k},
\]
e multiplicando por 2 a equação acima, obtemos
\[
	2 N = d_1 + d_2 \times \underbrace{2^{-1} + \dotsc + d_{k} 2^{-k+1}}_{N_1},
\]
ou seja, \(d_1\) é a ``parte inteira'' de \(2N\) e \(N_1\) é a ``parte fracionária''. Multiplicando \(N_1\) por 2 em ambos os lados, assim como antes, chegamos em
\[
	2 N_1 = d_2 + d_3 \times 2^{-1} + \dotsc + d_{k} \times 2^{-k+2},
\]
resultando na parte inteira de \((2 N_1)\) sendo \(d_2\) e, indutivamente, repete-se o processo até zerar a parte fracionária.
\begin{exr}
	Represente os números 0,125, 3,8 na base \(\beta =2\) [Spoiler: \(0,125 = (0,001)_{2}\) e \(3,8 = (11,110011001100 \dotsc )_{2}\)], e tente fazer um algoritmo para representar um número decimal fracionário em um número binário (o oposto também).
\end{exr}
\end{document}
